% !TEX root = ../../main.tex
\begin{tikzpicture}[
  font=\songti\zihao{5},
  node distance=1.0cm,
  startstop/.style={draw, rounded corners, align=center, minimum width=3.4cm, minimum height=0.82cm},
  process/.style={draw, rectangle, align=center, minimum width=5.0cm, minimum height=0.82cm},
  decision/.style={draw, diamond, aspect=2.2, align=center, inner sep=1pt, minimum width=4.8cm},
  arrow/.style={-{Stealth[length=2mm]}, thick}
]

\node[startstop] (start) {管理员在后台管理端选择模型版本并点击发布};
\node[process, below=of start] (request) {Spring Boot 服务端接收发布请求};
\node[decision, below=of request] (check) {模型文件与标签映射文件是否完整};
\node[process, below=of check] (reload) {服务端调用 Python 识别服务重载模型};
\node[decision, below=of reload] (success) {模型重载是否成功};
\node[startstop, below=of success] (end) {返回结果并结束处理};

\node[process, right=2.5cm of check] (fail1) {返回发布失败信息};
\node[process, right=2.5cm of success] (fail2) {返回模型加载失败信息};
\node[process, below=of fail2] (update) {更新数据库中的当前发布模型版本状态};
\node[process, below=of update] (record) {记录当前模型版本信息};
\node[process, below=of record] (result) {返回发布成功结果};
\node[process, below=of result] (usemodel) {移动端实时识别与视频识别使用当前发布模型};

\draw[arrow] (start) -- (request);
\draw[arrow] (request) -- (check);
\draw[arrow] (check) -- node[right] {是} (reload);
\draw[arrow] (check.east) -- ++(1.6,0) |- node[pos=0.20,above] {否} (fail1.north);
\draw[arrow] (fail1) |- (end.east);
\draw[arrow] (reload) -- (success);
\draw[arrow] (success) -- node[right] {成功} (update);
\draw[arrow] (success.east) -- ++(1.4,0) |- node[pos=0.20,above] {失败} (fail2.north);
\draw[arrow] (fail2) |- (end.east);
\draw[arrow] (update) -- (record);
\draw[arrow] (record) -- (result);
\draw[arrow] (result) -- (usemodel);
\draw[arrow] (usemodel) |- (end.west);

\end{tikzpicture}